
本课程设计以某三层办公楼为对象，完成供冷季制冷空调设计中的夏季冷负荷计算、房间级热源参数设置、新风处理和设备选型。建筑面积约1358.4\,m²，平面尺寸约32.0\,m $\times$ 14.15\,m，三层共划分为15个天正暖通计算分区，并在CLTD/CLF计算中保留50个房间的人员、照明、设备、新风和末端选型参数。设计选取沈阳、天津、成都、重庆和深圳五个城市进行比较，分析围护结构、室外温湿度和新风除湿需求对多联机和风机盘管加独立新风系统的影响。

全文采用天正暖通计算的工程设计思路：室内设计参数统一取夏季室温26\,\textcelsius、相对湿度60\%，基准新风量40\,m³/(h·人)，照明功率密度10\,W/m²，普通办公设备功率密度15\,W/m²；会议、门厅、医务、档案、卫生间和走廊等房间按功能分别修正人员密度、设备功率和新风指标。围护结构按五城市气候区差异化设置，且南方城市采用略优于参考设计的外窗遮阳和屋面保温参数，以体现本设计的节能优化取向。

计算结果表明，五城市三层办公楼设计峰值冷负荷由低到高大致为沈阳、成都、天津、深圳、重庆。重庆和深圳的新风及除湿负荷占比更高，设备选型受高温高湿和长供冷季共同约束；沈阳、天津则更强调围护保温与短时峰值冷量校核。多联机室内机按房间峰值冷负荷乘1.3安全系数选型，大房间采用多台嵌入式末端并联，每层设置独立室外机系统并校核室内外机容量配比。进一步设置新风量、室内温度、窗墙比、设备功率密度和围护遮阳的单因素灵敏度分析，用于识别影响冷源容量和新风负荷的关键参数。

\smallskip
\noindent\textbf{关键词：}制冷空调；办公建筑；冷负荷计算；供冷季；设备选型；多联机；风机盘管；独立新风
